We give an explicit example in which chronological order changes the analytic
type of a dynamical zeta factor.  Let the full two-shift drive two
noncommuting automorphisms of the dyadic two-solenoid
\(\widehat{\Z[1/2]^2}\), generated by
\(A=\left(\begin{smallmatrix}3&1\\1&3\end{smallmatrix}\right)\) and
\(B=\left(\begin{smallmatrix}3&2\\2&4\end{smallmatrix}\right)\).
For a chronological word \(w\), the fibre fixed group has the intrinsic
order
\[
 \#\Fix(\alpha_{M_w})
 =\oddpart\!\left(8^{|w|}-\tr M_w+1\right).
\]
The primitive words \texttt{aabbb} and \texttt{ababb} have the same letter
counts but fixed-point orders \(30035\) and \(15021\), respectively.  The
first return has a rational zeta function, whereas the second has a natural
boundary at base-clock radius \(1/8\).  Modulo two, the words carrying a
nontrivial local correction are exactly the cyclic golden-mean language.
Nevertheless, the full switching zeta does not inherit a natural boundary at
its first circle: it has a meromorphic, zero-free continuation apart from one
simple pole at \(1/16\) throughout
\(|z|<(8\golden)^{-1}\), where
\(\golden=(1+\sqrt5)/2\).  Thus orbit-resolved chronology and global
aggregation exhibit different analytic behavior.  The construction supplies
an exact arithmetic-dynamics example but fails the periodic-clock, divisor,
and operator requirements of a Hilbert--P\'olya model.
